% ============================================================================
% Check2HGI -- o andar de check-in (impresso ~30/31).
% Fonte: o ASCII do AUTOR, `considerations.md:3101-3130`.  Transcricao, nao
% reinterpretacao: Check-in -> Encoder -> Check-in Feature Embedding -> grafo
% de sequencia DIRIGIDA -> original/corrompido -> a MESMA GCN -> Check-in
% Embeddings -> Attention Pooling per POI -> POI Embeddings -> HGI.
%
% ESTA CHAPA E' IRMA DA DO HGI, DE PROPOSITO E NA LETRA.  Mesmos estilos, mesma
% paleta, mesmas raias, mesmos glifos de vetor, mesmo tratamento de "a mesma
% rede recebe os dois ramos".  E' o que torna literal o argumento do slide --
% "e' o mesmo metodo, um nivel mais fundo" -- em vez de pedir que a plateia
% acredite na palavra.  Por isso ela TERMINA numa caixa `HGI`: o que sai daqui
% e' o que entra na chapa anterior.
%
% ⚠ AS ARESTAS DO GRAFO SAO DIRIGIDAS, e isso nao e' decoracao.  A sequencia de
% visitas liga cada visita a's seguintes em UMA direcao so' (anterior -> posterior).
% E' o mecanismo antivazamento do v18: com a aresta bidirecional, o rotulo
% vazava e valia 28,63 macro-F1 em Alabama.  Desenhar sem ponta, ou com ponta
% nos dois lados, contradiz a defesa central do capitulo.  O proprio autor
% escreve "Build directed user visit sequence".
% ============================================================================
% O MESMO pictograma de grafo das chapas do DGI e do HGI -- mesmas posicoes de
% no', mesma silhueta -- so' que com as arestas DIRIGIDAS.  Duas razoes para nao
% escolher entre uma coisa e outra:
%  - consistencia: o leitor reconhece "isto e' um grafo" pelo mesmo simbolo nas
%    tres chapas, e a corrente de bolinhas em fila nao tinha esse parentesco;
%  - fidelidade: a sequencia de visitas e' DIRIGIDA e DESCONEXA, e as duas
%    propriedades sao substantivas:
%      * dirigida -- as arestas correm todas no sentido do tempo e NENHUMA fecha
%        ciclo.  Um ciclo dirigido diria que a ultima visita aponta de volta
%        para a primeira, que e' exatamente o que o v18 proibiu;
%      * desconexa -- a aresta de sucessao liga visitas DO MESMO USUARIO, entao
%        usuarios diferentes formam componentes separados.  Um grafo conexo
%        diria que todas as visitas se ligam entre si, o que e' falso.
%    E' aqui que esta chapa se separa das outras duas de proposito: o grafo de
%    Delaunay (DGI/HGI) e o de adjacencia de regioes (HGI) sao CONEXOS, e os
%    pictogramas deles refletem isso.  Mesmo simbolo, propriedade diferente.
\tikzset{
  dgpic/.pic={
    \begin{scope}[->, line width=0.8pt, opacity=0.55,
                  >={Latex[length=1.4mm]}, shorten >=1.4pt, shorten <=1.4pt]
      % Cada no' tem NO MAXIMO uma aresta de entrada e uma de saida -- ou seja,
      % cada componente e' um CAMINHO, nao um grafo qualquer.  E' o que torna a
      % coisa uma SEQUENCIA: uma visita tem uma anterior e uma seguinte, nunca
      % duas.  Ramificar diria que a mesma visita e' sucedida por duas outras.
      % Os nos ficam espalhados, e nao em fila, para o simbolo continuar sendo
      % o mesmo das chapas do DGI e do HGI -- muda a propriedade, nao o simbolo.
      %
      % Componente 1 -- as visitas de um usuario
      \draw (-0.54,-0.06) -- (-0.38, 0.28);
      \draw (-0.38, 0.28) -- (-0.06, 0.16);
      \draw (-0.06, 0.16) -- (-0.16,-0.22);
      % Componente 2 -- as visitas de OUTRO usuario, sem aresta ligando os dois
      \draw ( 0.22,-0.34) -- ( 0.54,-0.22);
      \draw ( 0.54,-0.22) -- ( 0.40, 0.06);
    \end{scope}
    \fill (-0.54,-0.06) circle (1.55pt) (-0.38,0.28) circle (1.55pt)
          (-0.06, 0.16) circle (1.55pt) (-0.16,-0.22) circle (1.55pt)
          ( 0.22,-0.34) circle (1.55pt) ( 0.54,-0.22) circle (1.55pt)
          ( 0.40, 0.06) circle (1.55pt);
  }
}
\begin{tikzpicture}[
    font=\small,
    >={Latex[length=2mm]},
    lbl/.style={align=center, font=\small, text=secondary},
    box/.style={draw=secondaryshade, rounded corners=2.5pt, fill=white,
                align=center, inner sep=4pt, line width=1pt, text=secondary},
    % caixas que atravessam as duas raias -- so' estas precisam ser altas
    tall/.style={box, minimum height=2.25cm},
    vcell/.style={draw=secondary!70, line width=0.8pt, minimum width=4mm,
                  minimum height=2.1mm, inner sep=0pt, fill=white},
    ncell/.style={vcell, draw=alert!85, fill=alert!10},
    flow/.style={->, line width=1.1pt, draw=secondary},
    nflow/.style={->, line width=1.1pt, draw=alert!85}
]
\def\yA{0.95}\def\yB{-0.95}

% O encoder e' de UMA raia so' -- nao precisa atravessar nada, e achatado ele
% para de competir com as caixas que de fato atravessam.
% ⚠ ERA "Encoder / pretrained", e a segunda palavra era FALSA (corrigido 27/08,
% medido pelo `gate`).  Nada pre-treinado esta' no caminho check-in -> feature:
% os 15 valores sao LIDOS do check-in (7 indicadores de categoria + 4 de tempo
% ciclico + 4 de tempo decorrido, `apx_h:95`, `:437`), e quem mapeia 15 -> 64 e'
% a PRIMEIRA CAMADA DE CONVOLUCAO do proprio Check2HGI (`apx_h:130-133`) -- que
% nesta chapa e' a caixa `GCN`, DEPOIS desta.  O que E' pre-treinado existe, mas
% e' outro andar: a ancora da place-table (`apx_h:167`), a jusante do pooling.
% `Preprocessing` e' tambem a primeira palavra do desenho do autor
% (`considerations.md:3128-3135`).  A caixa alargou 0,6cm e a PAGINA NAO MUDOU
% -- os extremos do desenho estao noutro lugar (verificado por export).
% ⚠ GEOMETRIA APERTADA -- medida, nao estimada.  O vao util vai da borda direita
% do rotulo "Check-in" (x=0,496; `Check-in` mede 1,191cm em \small) ate' a
% pilha `fe`.  E `fe` NAO pode ir muito para a direita: a seta seguinte faz
% cotovelo em x=4,08 e, se `fe.east`+0,22 passar disso, o trecho final INVERTE
% e a ponta aponta para tras (mesmo defeito que a chapa do DGI teve).
% Por isso: \footnotesize em vez de \normalsize, e a caixa recentrada em 1,80.
% ⚠ MEDIR NO BEAMER, nao em `article`.  A base do deck e' 11pt, entao \small=10pt
%   e \footnotesize=9pt -- um ponto acima do que article da'.  "Preprocessing"
%   mede 2,32cm a 11pt / 2,11cm a 10pt / 1,90cm a 9pt.  So' o ultimo cabe nos
%   1,95cm de `text width`.  Medir na classe errada custou uma iteracao aqui.
% Folgas resultantes: 0,19cm ate' o rotulo, 0,26cm ate' `fe`, 0,27cm de coto
% final da seta do cotovelo.  A PAGINA NAO MUDOU: 393,24 x 120,24 pt.
\node[box, minimum height=1.3cm, text width=1.95cm, inner sep=4pt] (enc) at (1.80,0)
      {\footnotesize\textbf{Preprocessing}};
\node[lbl, anchor=south] at (-0.10,0.10) {Check-in};
\draw[flow] (-0.30,0) -- (enc.west);

\begin{scope}[local bounding box=fe]
  \foreach \k in {0,1,2,3} \node[vcell] at (3.38,0.33-\k*0.22) {};\end{scope}
\draw[flow] (enc.east) -- (fe.west);

{\color{secondary}\pic at (4.75,\yA) {dgpic};}
{\color{alert}\pic at (4.75,\yB) {dgpic};}
\node[lbl, text width=1.9cm] at (4.90,0)
      {\textbf{directed visit}\\\textbf{sequence}};
\node[lbl, anchor=south] at (4.75,\yA+0.42) {original graph};
\node[lbl, text=alert, anchor=north] at (4.75,\yB-0.42) {corrupted graph};
\draw[flow]  (fe.east) -- ++(0.22,0) |- (4.08,\yA);
\draw[nflow] (fe.east) -- ++(0.22,0) |- (4.08,\yB);

\node[tall, text width=0.85cm] (gcn) at (6.70,0) {\textbf{GCN}};
\draw[flow]  (5.46,\yA) -- ([yshift=9.5mm]gcn.west);
\draw[nflow] (5.46,\yB) -- ([yshift=-9.5mm]gcn.west);

\begin{scope}[local bounding box=cea]
  \foreach \k in {0,1,2,3} \node[vcell] at (8.05,\yA+0.33-\k*0.22) {};\end{scope}
\begin{scope}[local bounding box=ceb]
  \foreach \k in {0,1,2,3} \node[ncell] at (8.05,\yB+0.33-\k*0.22) {};\end{scope}
\draw[flow]  ([yshift=9.5mm]gcn.east)  -- (cea.west);
\draw[nflow] ([yshift=-9.5mm]gcn.east) -- (ceb.west);

\node[tall, text width=1.6cm] (pool) at (9.50,0)
      {\textbf{Attention}\\\textbf{pooling}};
\draw[flow]  (cea.east) -- ([yshift=9.5mm]pool.west);
\draw[nflow] (ceb.east) -- ([yshift=-9.5mm]pool.west);

\begin{scope}[local bounding box=pea]
  \foreach \k in {0,1,2,3} \node[vcell] at (11.20,\yA+0.33-\k*0.22) {};\end{scope}
\begin{scope}[local bounding box=peb]
  \foreach \k in {0,1,2,3} \node[ncell] at (11.20,\yB+0.33-\k*0.22) {};\end{scope}
\draw[flow]  ([yshift=9.5mm]pool.east)  -- (pea.west);
\draw[nflow] ([yshift=-9.5mm]pool.east) -- (peb.west);
\node[lbl, text width=1.85cm, anchor=north] at (11.20,\yB-0.42) {place embeddings};

\node[box, fill=secondarytint, text width=0.75cm, minimum height=1.6cm]
     (hgi) at (12.45,0) {\textbf{HGI}};
\draw[flow]  (pea.east) -- ([yshift=5mm]hgi.west);
\draw[nflow] (peb.east) -- ([yshift=-5mm]hgi.west);

\end{tikzpicture}
